\chapter{Tinjauan Pustaka}

\section{Dasar Analisis Teknikal \& Representasi Data Finansial}
Analisis teknikal merupakan metodologi untuk memprediksi arah pergerakan harga di masa depan melalui studi data pasar historis.
\begin{itemize}
    \item \textbf{Visualisasi Runtun Waktu (Time-Series):} Representasi grafik harga di mana sumbu-X menunjukkan waktu ($t$) dan sumbu-Y menunjukkan harga ($P$). Fluktuasi harga dianggap sebagai pergerakan stokastik.
    \item \textbf{Anatomi Candlestick (OHLC):} Menggunakan empat variabel utama yaitu \textit{Open, High, Low,} dan \textit{Close} dalam satu jendela waktu ($\Delta t$).
    \item \textbf{Timeframe \& Sampling Data:} Transformasi data dari harga ($P_t$) menjadi \textit{Logarithmic Returns} ($R_t$) dilakukan untuk mencapai stasioneritas data.
\end{itemize}

\section{Dinamika Pasar \& Analisis Fundamental}
Analisis fundamental berfokus pada nilai intrinsik aset berdasarkan faktor ekonomi dan keuangan.
\begin{itemize}
    \item \textbf{Mekanisme Penggerak Harga:} Proses \textit{price discovery} melalui \textit{Auction Theory} (\textit{bid} dan \textit{ask}) serta hukum penawaran dan permintaan.
    \item \textbf{Metrik Valuasi \& Risiko:} Penggunaan rasio seperti P/E \textit{Ratio}, \textit{Dividend Yield}, serta pengukuran kinerja seperti \textit{Sharpe Ratio} dan \textit{Sortino Ratio}.
    \item \textbf{Makroekonomi \& Siklus Pasar:} Pemahaman terhadap rezim pasar (\textit{Bullish, Bearish, Sideways}) serta dampak kebijakan moneter (inflasi dan suku bunga).
    \item \textbf{Modern Portfolio Theory (MPT):} Konsep \textit{Efficient Frontier} untuk alokasi aset yang optimal secara klasik.
\end{itemize}

\section{Game Theory (Klasik)}
Teori permainan menyediakan kerangka matematis untuk menganalisis interaksi antar agen yang rasional.
\begin{itemize}
    \item \textbf{Model Matematika:} Melibatkan \textit{Payoff Matrix, Nash Equilibrium,} dan \textit{Pareto Optimum}.
    \item \textbf{Klasifikasi Permainan:} Terdiri dari \textit{Coordination, Anti-Coordination,} dan \textit{Mixed Strategy Game}.
    \item \textbf{Model Strategis Sekuensial:} Contohnya adalah \textit{Prisoner's Dilemma} sebagai analogi kegagalan pasar, dan model \textit{Stackelberg} untuk dinamika \textit{Leader-Follower}.
\end{itemize}

\section{Teori Informasi Kuantum \& Game Theory}
Penggabungan teori informasi kuantum ke dalam teori permainan memungkinkan ruang strategi yang lebih luas melalui superposisi dan keterkaitan (\textit{entanglement}).
\begin{itemize}
    \item \textbf{Kuantisasi EWL:} Skema Eisert-Wilkens-Lewenstein menggunakan operator \textit{entangling} ($\hat{J}$) dan strategi uniter untuk memetakan permainan klasik ke ranah kuantum.
    \item \textbf{Representasi State:} Menggunakan matriks densitas ($\rho$) dan pengukuran ketidakpastian melalui \textit{Von Neumann Entropy} ($S = -\text{Tr}(\rho \ln \rho)$).
    \item \textbf{Quantum Mutual Information (QMI):} Digunakan untuk mengukur korelasi total antar aset melalui pertukaran informasi kuantum.
\end{itemize}

\section{Quantum Algorithm (VQE)}
\textit{Variational Quantum Eigensolver} (VQE) adalah algoritma hibrid klasik-kuantum untuk mencari nilai eigen terendah dari sebuah Hamiltonian.
\begin{itemize}
    \item \textbf{Prinsip Variasi:} Menggunakan persamaan eigen $H|\psi\rangle = E|\psi\rangle$ untuk mencari \textit{ground state} yang diinterpretasikan sebagai kondisi ekuilibrium pasar paling stabil.
    \item \textbf{Struktur Ansatz:} Terdiri dari rotasi qubit tunggal ($R_y, R_z$) untuk pemetaan ruang strategi dan \textit{layer entanglement} (CNOT) untuk memodelkan interdependensi aset.
    \item \textbf{Teknik Optimasi:} Menggunakan optimisator klasik seperti SPSA (\textit{Simultaneous Perturbation Stochastic Approximation}) atau COBYLA.
    \item \textbf{Ising Hamiltonian:} Pemetaan masalah optimasi ke dalam model Ising dengan parameter bias ($h_i$) dan interaksi ($J_{ij}$).
\end{itemize}

\newpage
\thispagestyle{empty}
\vspace*{\fill}
    \begin{center}
	Halaman ini sengaja dikosongkan
    \end{center}
\vspace*{\fill}
\pagebreak
